\section{软件设计与流程}

\subsection{软件总体架构}

系统软件采用分层状态机架构，自下而上包括驱动层、采样与保护层、控制算法层、能量管理层和通信交互层。驱动层负责 ADC、PWM、定时器、通信接口和 GPIO 初始化；采样与保护层完成数据滤波、量纲转换、阈值判断和故障锁存；控制算法层运行 MPPT、母线稳压、充放电电流控制和交流输出控制；能量管理层根据光伏功率、负载功率、SOC 和故障状态选择工作模式；通信交互层完成 MCU 间数据同步和无线监控。



\begin{figure}[htbp]
  \centering
  \resizebox{0.96\linewidth}{!}{
  \begin{tikzpicture}[
    font=\small,
    node distance=1.0cm and 1.2cm,
    process/.style={
      draw,
      rounded corners=2pt,
      align=center,
      minimum width=2.8cm,
      minimum height=0.9cm,
      inner sep=4pt
    },
    line/.style={-Latex, thick}
  ]

  % ===================== 主流程 =====================
  \node[process] (poweron) {系统上电};
  \node[process, right=of poweron] (init) {外设初始化};
  \node[process, right=of init] (selfcheck) {系统自检};
  \node[process, right=of selfcheck] (precharge) {预充 / 软启动};
  \node[process, right=of precharge] (mode) {模式判断};

  % ===================== 运行流程 =====================
  \node[process, below=1.5cm of mode] (control) {闭环控制};
  \node[process, left=of control] (upload) {数据上报\\故障检测};
  \node[process, left=of upload] (fault) {故障判断};

  % ===================== 结果状态 =====================
  \node[process, below=1.5cm of upload] (run) {正常运行};
  \node[process, below=1.5cm of fault] (protect) {保护停机};

  % ===================== 主流程箭头 =====================
  \draw[line] (poweron) -- (init);
  \draw[line] (init) -- (selfcheck);
  \draw[line] (selfcheck) -- (precharge);
  \draw[line] (precharge) -- (mode);

  % ===================== 进入闭环 =====================
  \draw[line] (mode.south) -- node[left]{满足运行条件} (control.north);

  % ===================== 闭环与检测 =====================
  \draw[line] (control.west) -- (upload.east);
  \draw[line] (upload.west) -- (fault.east);

  % ===================== 故障判断 =====================
  \draw[line] (fault.south) -- node[above]{否} (run.north);
  \draw[line] (fault.south) -- node[left]{是} (protect.north);

  % ===================== 正常运行周期循环：走外圈 =====================
  \draw[line] (run.east) -- ++(5,0) |- node[above]{周期循环} (mode.east);

  % ===================== 故障恢复：走下方外圈 =====================
  \draw[line] (protect.west) -- ++(-0.75,0) |- node[pos=0.22, left]{故障清除 / 重新自检} (selfcheck.west);

  \end{tikzpicture}
  }
  \caption{系统主程序流程图}
  \label{fig:main_program_flow}
\end{figure}



\subsection{前级三端口控制策略}

前级三端口控制的核心目标是协调光伏、电池和直流母线。系统根据光伏功率 $P_{\mathrm{pv}}$、负载功率 $P_{\mathrm{load}}$、电池 SOC 和故障状态划分工作模式：

\begin{enumerate}[label=\arabic*、, leftmargin=*, itemsep=0.8ex, parsep=0pt, topsep=0.5ex]
  \item 光伏独立供电模式：当光伏功率接近负载需求且电池无需充放电时，三端口变换器优先维持母线稳定。
  \item 光伏供电并充电模式：当光伏功率大于负载需求且电池允许充电时，多余能量进入电池，充电电流受 BMS 限值约束。
  \item 光伏与电池联合供电模式：当光伏功率小于负载需求且电池允许放电时，电池端口补充能量以维持直流母线。
  \item 电池离网供电模式：当光伏输入不足或夜间运行时，电池作为主要能量来源，系统根据 SOC 限制输出功率。
  \item 保护降额或停机模式：当出现过压、欠压、过流、过温或通信异常时，系统进入限功率或停机状态。
\end{enumerate}

MPPT 算法可采用扰动观察法作为基础实现。控制器周期性调整光伏端等效工作点，比较扰动前后的光伏功率变化，决定下一周期扰动方向。为避免快速负载扰动导致误判，MPPT 更新频率应低于电流环和母线环，同时在母线异常或电池保护状态下暂停最大功率跟踪。

\subsection{后级交流变换器控制策略}

后级交流变换器以输出电压稳定和频率准确为主要目标。离网模式下，交流变换器生成 50 Hz 正弦参考，通过电压外环和电流内环控制全桥 PWM，使输出电压满足交流负载供电需求。输出 LC 滤波器参数需要与控制器带宽匹配，避免轻载和非线性负载下产生谐振。

交流控制与前级母线控制之间需要协调。当负载阶跃导致母线电压下降时，后级控制应限制瞬时电流指令，前级控制提高供能功率；当母线过压时，交流变换器可提高输出消纳能力或通知前级降低光伏输入功率。该协调逻辑可以减少母线过冲和输出波形畸变。

\subsection{多 MCU 协同控制策略}

三颗 MCU 之间采用主从协同关系。GD32G553 是功率控制主节点，周期性接收 BMS 状态和无线配置参数，并向无线模块发送运行数据。GD32C113 是电池安全节点，向主控提供 SOC、允许充电电流、允许放电电流和故障状态。GD32VW553 是通信节点，负责将用户配置转换为主控可识别的参数命令。

通信协议建议采用固定帧格式，包括帧头、设备地址、命令字、数据长度、数据区和校验字段。关键控制参数需要设置范围校验和生效条件，避免无线侧误操作导致系统异常。若主控在设定时间内未收到 BMS 心跳，应进入电池通信异常状态，并限制充放电功率或停机保护。

\begin{table}[htbp]
  \centering
  \caption{MCU 间关键交互数据}
  \label{tab:mcu_data}
  \begin{tabular}{|p{0.18\linewidth}|p{0.26\linewidth}|p{0.42\linewidth}|}
    \hline
    数据方向 & 数据内容 & 用途 \\
    \hline
    BMS 到主控 & SOC、单体极值、温度、充放电允许电流、故障码 & 决定能量调度、限流和保护状态 \\
    \hline
    主控到 BMS & 工作模式、端口电流、电池功率、系统故障状态 & 辅助 BMS 判断系统级风险 \\
    \hline
    主控到无线模块 & 电压、电流、功率、模式、故障码 & 状态显示、日志记录和远程监控 \\
    \hline
    无线模块到主控 & 目标模式、限值参数、复位/清故障命令 & 用户配置和运维控制 \\
    \hline
  \end{tabular}
\end{table}

\subsection{上位机与无线监控设计}

无线监控模块主要完成数据可视化、参数配置和故障告警。监控界面建议展示光伏电压电流、光伏功率、电池电压电流、SOC、直流母线电压、交流输出电压、输出功率、系统模式和故障码。参数配置功能包括 MPPT 使能、输出开关、充放电电流限值、SOC 上下限和日志导出。

\todo{补充上位机界面截图、通信协议字段定义和实际无线链路测试结果。}
